\documentclass[12pt,a4paper]{report}
\usepackage[utf8]{inputenc}
\usepackage{amsmath}
\usepackage{amsfonts}
\usepackage{amssymb}
\usepackage{amsthm}
\usepackage{hyperref}

\usepackage{multicol}
\usepackage{fancyhdr}
\usepackage[inline]{enumitem}
\usepackage{tikz}
\usepackage{tikz-cd}
\usetikzlibrary{calc}
\usetikzlibrary{shapes.geometric}
\usetikzlibrary{positioning}
\usepackage[margin=0.5in]{geometry}
\usepackage{xcolor}

\hypersetup{
    colorlinks=true,
    linkcolor=blue,
    filecolor=magenta,      
    urlcolor=cyan,
    pdftitle={Tensors},
    pdfpagemode=FullScreen,
    }

%\urlstyle{same}

\newcommand{\CLASSNAME}{Math 5050 -- Special Topics: Tensor Analysis}
\newcommand{\STUDENTNAME}{Paul Carmody}
\newcommand{\ASSIGNMENT}{Notes and Observations }
\newcommand{\DUEDATE}{January 2026}
\newcommand{\PROFESSOR}{Professor Berchenko-Kogan}
\newcommand{\SEMESTER}{Spring 2026}
\newcommand{\SCHEDULE}{TBD}
\newcommand{\ROOM}{Remote}

\newcommand{\MMN}{M_{m\times n}}
\newcommand{\FF}{\mathcal{F}}

\pagestyle{fancy}
\fancyhf{}
\chead{ \fancyplain{}{\CLASSNAME} }
%\chead{ \fancyplain{}{\STUDENTNAME} }
\rhead{\thepage}

\fancyfoot{} % clear all footer fields
\fancyfoot[LE,RO]{\thepage}
\fancyfoot[LO,CE]{\ASSIGNMENT}
%\fancyfoot[CO,RE]{To: Dean A. Smith}
\input{../MyMacros}

\newcommand{\RED}[1]{\textcolor{red}{#1}}
\newcommand{\BLUE}[1]{\textcolor{blue}{#1}}

\newcommand{\SUPP}{\operatorname{supp}}
\newcommand{\CLOSURE}{\operatorname{closure of}}
\newcommand{\BD}{\operatorname{bd}}
\newcommand{\HHN}{\mathcal{H}^n}
\newcommand{\COFF}[3]{\Gamma^{#1}_{{#2}{#3}}}
\newcommand{\VK}[1]{\textbf{#1}}

\begin{document}

\begin{center}
	\Large{\CLASSNAME -- \SEMESTER} \\
	\large{ w/\PROFESSOR}
\end{center}

\begin{center}
\textbf{\Large{Notes and Observations}}
\end{center}

\HLINE

\begin{description}

	\item \textbf{Semester Plan}

		One chapter per week with the hopes of all 17 chapters by the end of April.	
	
	\item \textbf{Status:}
	\begin{description}
		\item \textbf{January 20, 2026}  Reading Chapter 8 and Chapter 7 exercises.
	\end{description}

	\item \textbf{Important Quotes}
	
	``A tensor is not meaningful by itself $\dots$ tensors lead to invariants." -- Page 86
	
	\item \textbf{Indices:}
	
	The upper index is for covariant tensors and the lower index is for contravariant tensors.  Notice that 
	\begin{align*}
		\PART{f^i}{x^j}
	\end{align*}that $f$ is a function and $x$ is a vector.  This is in keeping with the definition from \textit{Manifolds} where a covariant tensor is a functional and a contravariant tensor is vector.  In the special case of the dual space then we know that a basis vector is of the form $f^i(x_j)=\delta^i_j$ which derives like a constant.  In less precise terms, a functional is made up of terms based on the coordinates.  Extending that notion, a tensor is rearrangement of coordinates (and applied in an expression) that responds to a change of coordinates in a specific way.
	
	\item \textbf{Progression of Intrinsic Objects}
	\begin{itemize}
		\item Start with scalar field.
		\item Given a Position Vector $\VK{R}(Z^1, Z^2, Z^3$) (no coordinates).
		\item Default tangent space $\VK{Z}_i = \PART{\VK{R}}{Z^i}$ also called  covariant tensors.
		\begin{enumerate}
			\item \textbf{Vector notation}
			
			\begin{align*}
				\VK{V} &= \sum_{i=1}^n V^i\VK{Z}_i = V^i\VK{Z}_i
			\end{align*}It seems a little silly but this $V^i\VK{Z}_i$ is a new intuitive form of a vector (the index is arbitrary and unique).
			
		\end{enumerate}
		
		\item Covariant metric tensor $Z_{ij} = \VK{Z}_i\cdot\VK{Z}_j$.  The foundational object of Reimannian Geometry.
		\item Contravariant metric tensor $Z^{ij}Z_{ij} = \mathbb{I}$ (inverse of the covariant metric tensor).
		\item Contravariant tensor derived from contravariant metric tensor, $\VK{Z}^j = Z^{ij}\VK{Z}_i$
	\end{itemize}
	
	\item Useful identitites
	\begin{align*}
		Z_{ij}Z^{ij} &= \mathbb{I} \\
		Z_{ij}Z^{ik} &= \delta^k_j
	\end{align*}
	
	\item a little dissection*
	\begin{align*}
		L &= \int^b_a \sqrt{Z_{ij}\frac{dZ^i}{dt}\frac{dZ^j}{dt}}\, dt \\
		&= \int^b_a \sqrt{\sum_{i=1}^n\sum_{j=1}^nZ_{ij}\frac{dZ^i}{dt}\frac{dZ^j}{dt}}\, dt \\
		&= \int^b_a \sqrt{\sum_{i=1}^n\sum_{j=1}^n\PAREN{\VK{Z}_i\cdot\VK{Z}_j}\frac{dZ^i}{dt}\frac{dZ^j}{dt}}\, dt & \VK{Z}_i\perp\VK{Z}_j, i \ne j\\
		&= \int^b_a \sqrt{\sum_{i=1}^n\PAREN{\frac{dZ^i}{dt}}^2}\, dt
	\end{align*}
	\item \textbf{Christoffel Symbols}
	
	By definitions:
	\begin{align*}
		\COFF{k}{i}{j} &= \VK{Z}^k\cdot \PART{\VK{Z}_i}{Z^j} & (5.60)\\
		&= -\PART{\VK{Z}^k}{Z^j}\cdot\VK{Z}_i & (5.64)
	\end{align*}Note: this is NOT a contraction.  Picture three $3\times 3$ matrices with these scalars.  Each element is the projection of a partial derivative  of a basis vector ($i$ with respect to $j$) onto another basis vector, $k$.  The collection of the elements of $k$ make up the coordinates of the partial $i$ with respect to $j$.\\
	
	The contravariant version is 
	\begin{align*}
		\PART{\VK{Z}^k}{Z^j} &= -\COFF{k}{i}{j}\VK{Z}^i & (5.65)
	\end{align*}note the contraction on $i$.\\
	
	And based on the metric tensors
	\begin{align*}
		\COFF{k}{i}{j} &= \frac{1}{2}Z^{km}\PAREN{\PART{Z_{mi}}{Z_j}+\PART{Z_{mj}}{Z_i}-\PART{Z_{ij}}{Z_m}} & (5.67)
	\end{align*}
	
	\textbf{Transformation of Coordinates}
	\begin{align*}
		\COFF{i}{j}{k} &= \CONJ{\Gamma}^u_{rs} \PART{x^i}{\CONJ{x}^u}\PART{\CONJ{x}^r}{x^j}\PART{\CONJ{x}^s}{x^k}+\frac{\partial^2\CONJ{x}^r}{\partial x^j\partial x^k}\PART{x^i}{\CONJ{x}^r}
	\end{align*}
	
	\item \textbf{Multiple change of coordinates}
	
	I find the notation for the Jacobian misleading at times.  We must understand that $J^i_j$ does not mean indices of $i$ and $j$ but different change of coordinate systems.  Further the over use of $Z^j$ when we might be meaning an entirely different basis vector gets confusing as well.  Suppose $a$ represents Cartesian Coordinates $(x,y)$ and $b$ represents polar coordinates $(r, \theta)$ then
	\begin{align*}
		J^a_b &= \TWOXTWO{x_r}{x_\theta}{y_r}{y_\theta}
	\end{align*}using the susbscript partial derivative notation.  Now, let $c$ represent an affine coordinate system $p,q$  then 
	\begin{align*}
		J^b_c &= \TWOXTWO{r_p}{r_q}{\theta_p}{\theta_q} 
	\end{align*}Notice that combining these two we get
	\begin{align*}
		J^a_bJ^b_c &= \TWOXTWO{x_r}{x_\theta}{y_r}{y_\theta}\TWOXTWO{r_p}{r_q}{\theta_p}{\theta_q} \\
		&= \TWOXTWO{x_rr_p+x_\theta\theta_p}{x_rr_q+x_r\theta_q}{y_rr_p+y_\theta\theta_p}{y_rr_q+y_r\theta_q}
	\end{align*}Notice that $x_rr_p = \PART{x}{r}\PART{x}{p}= \PART{x}{p}$.  Thus,
	\begin{align*}
	J^a_bJ^b_c &= \TWOXTWO{x_p+x_p}{x_q+x_q}{y_p+y+p}{y_q+y_q} = 2J^a_c
	\end{align*}the $b$ in the superscript followed by the subscript of $J$ effectively cancel out, but $b$ is not an index.  Remember that it represents a change of coordinates that has become superfluous.
	
	\item \textbf{On Page 76} We have the identities (4.74) and (4.76)
	\begin{align*}
		J^i_{i'}J^{i'}_j &= \delta^i_j & (4.74)\\
		J^{i'}_iJ^i_{j'} &= \delta^{i'}_{j'} & (4.76)
	\end{align*}\BLUE{Abuse of notation.}  From above we have $J^a_bJ^b_c = J^a_c$ where each $J$ is a distinct coordinate transformation between three different systems.  Here, though, we are actually saying that $i,j$ and $i',j'$ are the same coordinate system but with different indices.  We can see the contradiction in $J^a_c$.  This implies $\delta^a_c = 1$ if $x=p$ (which might be true in the affine case).  Consider in stead
	\begin{align*}
		J^b_a &= \TWOXTWO{r_x}{r_y}{\theta_x}{\theta_y} \\
		J^b_a J^a_c &= \TWOXTWO{r_xx_p+r_yy_p}{r_xx_q+r_yy_q}{\theta_xx_p+\theta_yy_p}{\theta_xx_q+\theta_yy_q} \\
		&= 2\TWOXTWO{r_p}{r_q}{\theta_p}{\theta_q}
	\end{align*}which is distinctly NOT $\delta^b_c$.
	
	\item \textbf{On page 80} we see this abuse again where $i,j$ and $i',j'$ are considered different indices of the same coordinate transformation.
	
	\item \textbf{Beware the metric tensors.}  Notice that the covariant metric tensor is defined as
	\begin{align*}
		Z_{ij} &= \VK{Z}_i\cdot\VK{Z}_j.
	\end{align*}This only makes sense when $i,j$ are indices of the same coordinate system.  The notation, though, suggests that they could be from different coordinate systems.  Which might look something like
	\begin{align*}
		Z_{y\theta} &= y\cdot \theta &(?)
	\end{align*}Indeed, later expressions treat them that way as in 
	\begin{align*}
		Z_{ij}J^i_{i'}J^j_{j'} &= Z_{i'j'}
	\end{align*}In this case, $J^i_{i'}$ and $J^j_{j'}$ are the same matrix (and must be the same) with different indices.
	
	\item \textbf{Page 85} we say ``$S^i_j$ is a tensor then $S^i_i$ is an invariant."  In this case, can be different coordinate systems.  Let $i$ be Cartesian and $j$ be polar coordinates.  $J^j_i$ would be a change of coordinates from polar to Cartesian which makes $\bar{S}^i_i=S^i_jJ^j_i$ which would make sense but $\bar{S}$ is not made up of the same coordinate systems as $S$.  Be wary of $S^i_i$ when the subscript is originally from a different coordinate system.  The $S$ refers to the tensor object $S$ but $S^i_i$ and $S^i_j$ are very different flavors of it.  This can be deceptive.
	
	\item \textbf{Covariant Derivative $\nabla_j$}
	\begin{itemize}
		\item Contravariant Tensor
		
		\begin{align*}
			\nabla_j V^i = \PART{V^i}{Z^j}+\COFF{i}{j}{m}V^m
		\end{align*}
		
		\item Covariant Tensor
		
		\begin{align*}
			\nabla_j V_i = \PART{V_i}{Z^j}-\COFF{m}{i}{j}V_m
		\end{align*}

	\end{itemize}
	
	\item \textbf{Index Algebra}
	\begin{enumerate}
		\item The Free Index.  Indices appearing once on one side must appear once on ther other side.
		\item Indices appearing twice on one side must appear in both upper and lower positions.  Appearing on both sides actually means that the index on the left and the index on the right are independent of each other (this should be avoided).
		\item \textbf{Moving from left to right: }An element/tensor (i.e., a factor in multiplication) transfers from one side to the other (i.e., analogous to multiplying by its inverse) by introducing a free index and switching positions yielding (that is, $J_b^k \to J^b_m$ or $Z_{im} \to Z^{mj}$ because $J_b^kJ^b_m = \delta^k_m$ and $Z_{im}Z^{mj} = \delta_i^j$)\footnote{An expression that results in a Kronecker Delta is being considered a quotient expression}. For example,
		\begin{align*}
			J_b^k = \delta_b^c &\implies J_b^k J^b_m = \delta_b^c J^b_m \implies \delta_m^k= \delta_b^c J^b_m \implies \delta_m^k = J_m^c & (1)\\
			Z_{im} = J^c_m &\implies Z_{im}Z^{mj} = J^c_mZ^{mj} \implies \delta_i^j = J^c_mZ^{mj} & (2)
		\end{align*}(1) moves the $\delta$ from one side to the other side by removing $b$ and introducing $m$ (both free indices).  (2) moves $Z$ to the right (leaving the $\delta$ behind) by removing $i$ and introduction $j$ both of which are free indices.
	\end{enumerate}
	
	\item \textbf{Delta symbols}
	
	This symbol determines/filters two permuations that are in sync, out of sync, or neither.
	\begin{align*}
		e^{ijk} &= e_{ijk} = \BINDEF{1 & ijk \text{ positive permutation}}{-1 & ijk \text{ negative permuation} \\ 0 & \text{otherwise}} \\
		\delta_{rst}^{ijk} &= e_{rst}e^{ijk}
	\end{align*}Thus, $\delta_{rst}^{ijk}$ is positive when $ijk$ and $rst$ have the same permutation sign, negative when they have different permutation signs and zero otherwise.  Also, the numbers $ijk$ and $rst$ must come from the same set albeit not the same order (e.g., $\exists i \in rst$ as well as $j$ and $k$)
	
	Interesting identities
	\begin{align*}
		\delta_{rs}^{ij} &= \delta_r^i\delta_s^j - \delta_r^j\delta_s^i \\
		&= \left | \begin{array}{cc}
			\delta_r^i & \delta_s^i \\
			\delta_r^j & \delta_s^j
		\end{array} \right |
	\end{align*}
	
	\HLINE
	\section*{Chapter 10 -- Surface Tensors}
	\item \textbf{Shift Tensor}
	
	The matrix which converts from the ambient basis, $\textbf{Z}_i$, and the surface basis, $S_\alpha$.
	\begin{align*}
		Z_\alpha^i &= \textbf{S}_\alpha \cdot \textbf{Z}^i
	\end{align*}for example. 
	\begin{align*}
		\LET \textbf{T} &= T^i\textbf{Z}_i \\
			&= T^\alpha\textbf{S}_\alpha = T^\alpha Z_\alpha^i\textbf{Z}_i \\
			\therefore T^i &= T^\alpha Z_\alpha^i
	\end{align*}
	
	\item \textbf{Surface covariant metric tensor}
	
	Analogous to the covariant metric tensor $Z_{ij}$.
	\begin{align*}
		S_{\alpha\beta} &= \textbf{S}_\alpha \cdot \textbf{S}_\beta
	\end{align*}Consequently, the covariant metric tensor and the surface covariant metric tensor are related by 
	\begin{align*}
		S_{\alpha\beta} = \textbf{S}_\alpha\cdot\textbf{S}_\beta = \textbf{Z}_i Z_\alpha^i\,\cdot\,\textbf{Z}_j Z_\beta^i = Z_{ij}Z_\alpha^i Z_\beta^j
	\end{align*}The last term is all matrix multiplication starting with the metric tensor that twice the surface metric tensor.
	
	\item \textbf{Surface Contravariant Metric Tensor}.
	
	This is defined as the matrix inverse that provides
	\begin{align*}
		S^{\alpha\beta}S_{\beta\gamma} &= \delta_\gamma^\beta
	\end{align*}\textbf{Note:} $Z_{ij}, Z^{ij}$ would be a 3 dimensional square matrix, $S_{\alpha\beta}, S^{\alpha\beta}$ is a two dimensional matrix, $Z_\alpha^i$ is $2 \times 3$ matrix and $Z_i^\alpha$ is a $3 \times 2$ matrix.
	
	\item \textbf{Surface versions of Levi-Civita symbols}
	
	Define $S$ (not a vector and no subscripts) as anlogous ot $Z$ (not a vector and no subscripts) as
	\begin{align*}
		S = |S_{..}|
	\end{align*}We define the \textbf{area element} as $\sqrt{S}$.
	
	Define the permutation operator as 
	\begin{align*}
		e_{12} = -e_{21} = 1 \text{ zero otherwise}
	\end{align*}then
	\begin{align*}
		\varepsilon_{\alpha\beta} &= \sqrt{S}\,e_{\alpha\beta} \\
		\varepsilon^{\alpha\beta} &= \frac{e^{\alpha\beta}}{\sqrt{S}} 
	\end{align*}

\HLINE

\newpage
\item \textbf{Commutativity by Contraction.}

This is a really central structural fact about tensors, and it’s worth seeing both \DEFINE{how it works algebraically} and \DEFINE{why it must be true geometrically}.

---

\begin{description}
\item \textbf{1. The statement (what we want to understand)}

Let (T) be a tensor. Then:

\begin{align*}
\boxed{
\nabla_\mu(\text{contraction of } T) = \text{contraction of } (\nabla_\mu T) }
\end{align*}

In words:

\textbf{ \DEFINE{Covariant differentiation commutes with contraction}}

---

\item \textbf{2. Start with a concrete example}

Take a $(1,1)$-tensor:
\begin{align*}
T^\alpha_{\ \beta}
\end{align*}

Its contraction is:
\begin{align*}
T^\alpha_{\ \alpha}
\end{align*}

Now compute both sides.

---

\begin{description}
\item (A) Differentiate after contraction

\begin{align*}
\nabla_\mu (T^\alpha_{\ \alpha}) = \partial_\mu (T^\alpha_{\ \alpha})
\end{align*}

(since this is now a scalar)

---

\item (B) Contract after differentiating

\begin{align*}
\nabla_\mu T^\alpha_{\ \beta} = \partial_\mu T^\alpha_{\ \beta} + \Gamma^\alpha_{\mu\nu} T^\nu_{\ \beta} - \Gamma^\nu_{\mu\beta} T^\alpha_{\ \nu}
\end{align*}

Now contract $\alpha = \beta$:

\begin{align*}
\nabla_\mu T^\alpha_{\ \alpha} = \partial_\mu T^\alpha_{\ \alpha} + \Gamma^\alpha_{\mu\nu} T^\nu_{\ \alpha} - \Gamma^\nu_{\mu\alpha} T^\alpha_{\ \nu}
\end{align*}

---
\end{description}

\item \textbf{3. The key cancellation}

Look at the last two terms:

\begin{align*}
\Gamma^\alpha_{\mu\nu} T^\nu_{\ \alpha} \quad \text{and} \quad \Gamma^\nu_{\mu\alpha} T^\alpha_{\ \nu}
\end{align*}

Rename dummy indices in the second term:
\begin{align*}
\nu \leftrightarrow \alpha
\end{align*}

It becomes:
\begin{align*}
\Gamma^\alpha_{\mu\nu} T^\nu_{\ \alpha}
\end{align*}

\textbf{ So the two terms are identical and cancel:}

\begin{align*}
+ (\cdot) - (\cdot) = 0
\end{align*}


---

\item \textbf{4. Result}

\begin{align*}
\boxed{
\nabla_\mu (T^\alpha_{\ \alpha}) = \partial_\mu (\textbf{T}^\alpha_{\ \alpha}) = (\nabla_\mu T^\alpha_{\ \beta})\big|_{\alpha=\beta} 
}
\end{align*}

So contraction and covariant derivative commute.

---

\item \textbf{5. Why this always works}

Every contraction pairs:

\begin{itemize}
\item one \DEFINE{upper index}
\item one \DEFINE{lower index}
\end{itemize}

And in the covariant derivative:

\begin{itemize}
\item upper index contributes a $+\Gamma$ term
\item lower index contributes a $-\Gamma$ term
\end{itemize}

\textbf{ These always cancel when the indices are the same.}

---

\item \textbf{6. General principle}

\begin{align*}
\boxed{
\text{Each contraction produces a $+\Gamma$ and $-\Gamma$ pair that cancels}
}
\end{align*}

So this works for \DEFINE{any tensor}, not just $(1,1)$.

---

\item \textbf{7. Deep geometric reason}

Contraction is just:
\begin{align*}
\text{pairing } V^\alpha \text{ with } \omega_\alpha
\end{align*}

This pairing is \DEFINE{coordinate-invariant}.

The covariant derivative is built to preserve:

\begin{itemize}
\item tensor type
\item tensor operations
\end{itemize}

So it must satisfy:
\begin{align*}
\nabla(\text{pairing}) = \text{pairing of } \nabla
\end{align*}

---

\item \textbf{8. Even deeper: compatibility with the metric}

Because:
\begin{align*}
\nabla_\mu g_{\alpha\beta} = 0
\end{align*}

raising/lowering indices also commutes with $\nabla$, which is closely related to contraction.

---

\item \textbf{9. Clean abstract statement}

\begin{align*}
\boxed{
\nabla \circ \text{contraction} = \text{contraction} \circ \nabla }
\end{align*}

---

\item \textbf{10. Intuition}

\begin{itemize}
\item Contraction = “collapse two indices into a number”
\item Covariant derivative = “differentiate consistently”
\end{itemize}

\textbf{ Doing these in either order gives the same result because the connection treats upper/lower indices in exactly opposite ways.}

---

\item \textbf{11. One-line takeaway}

\begin{align*}
\boxed{
\text{Covariant differentiation respects tensor algebra operations.}
}
\end{align*}

---
\end{description}

If you want, I can next:

\begin{itemize}
\item show how this property \DEFINE{fails} for ordinary partial derivatives (nice contrast), or
\item connect this to why divergence has that clean $\frac{1}{\sqrt{|g|}}\partial(\sqrt{|g|}\cdot)$ form.
\end{itemize}


\newpage
\section*{Chapter 11 and 12 -- Curvature}

\item \textbf{Riemann-Christoffel Tensor}

\item \textbf{Curvature Tensor $B_{\alpha\beta}$ }

\begin{itemize}
	\item Definition (covariant form): $B_{\alpha\beta}$ is defined as the ``array" (not a matrix) of change in the shift tensor with respect to the a
		\begin{align*}
			B_{\alpha\beta} = \textbf{N} \nabla_\alpha S_{ \beta}
		\end{align*}
		$B_{\alpha\beta}$ measures how the surface bends in the normal direction.  It takes two tangent directions and gives a scalar.

	\item Invariant form $B^\alpha_\beta$: Raise the index using the surface metric tensor producing a matrix:
		\begin{align*}
			B^\alpha_{\beta} = g^{\alpha\gamma} B_{\gamma\beta}
		\end{align*}
		This is used as a linear operator transforming a vector to another.  Known as the \DEFINE{shape operator or Weingarten map}.

	\item Principle Curvature (eigenvalues of $B^\alpha_\beta$):
	\begin{align*}
		B^\alpha_\beta v^\beta = \kappa v^\alpha
	\end{align*}Where $\kappa$ is an eigenvalue of $B^\alpha\beta$ (one of two).  Raising the index of $B_{\alpha\beta} \to B^\alpha\beta$ allows for diagonalization.

	\item Mean curvature $B^\alpha_\alpha$ and trace $|B^._.|$
		\begin{align*}
			B^\alpha_\alpha &= \kappa_1 + \kappa_2 \\
			|B^._.| &= \det B^\alpha_\alpha = \kappa_1 \kappa_2
		\end{align*}

	\item \textbf{Primary point}.  The $B$ tensor is an example of a single object with many perspectives or uses depending both on the coordinates used and on the application (scalar curvature or linear operator, etc.)

	\item \textbf{Matrix Square and Groundforms}
		\begin{align*}
			C^\alpha_\beta &= B^\alpha_\gamma B^\gamma_\beta.
		\end{align*}Defines the following groundforms
		\begin{itemize}
			\item $1\FST$ ground form $S_{\alpha\beta}$ the metric tensor.
			\item $2\SND$ ground form $B_{\alpha\beta}$ the curvature tensor.
			\item $3\TRD$ ground form $C^\alpha_\beta$ the matrix square of the curvature tensor.
		\end{itemize}
\end{itemize}

\newpage
\section*{Interesting Notes}
\item \textbf{Interesting notes}
\begin{description}
	\item \textbf{What is a vector?}

		Think of it first as an object that has more than just a scalar value (e.g., temperature, humidity or luminosity) but also has direction (e.g., velocity, acceleration, gradiant of a scalar -- direction of greatest increase). These are depicted as arrows fixed at a point, with a length and pointing in a single direction.  By seeing it as the last example, direction of a scalar with greatest increase, we have two very important perspectives.  These will be in reverse order of the traditional introduction to vectors.
		\begin{itemize}
			\item \textbf{The perspective of the Covector.}

			This is really quite simple.  Picture the contours of the scalar values over the space (usually a surface).  Each contour is a line representing a constant value of our scalar (in a topographical map these lines would represent altitude).  The value of the covector \textit{of this object} at a point is the number of contours that the object (in the form of a line) passes through.  Thus, by definition the covector describes the space in question by describing how any of these objects passes.

			\item \textbf{The perspective of the Vector.}

			This aspect describes the position, length and direction of the object defined by coordinates associated with the contours (i.e., scalar field).  The position is deterined by bases that define the ambient space that the scalar field is in.  The direction is determined through the gradiant of the scalar field (directon of greatest change) and directional derivatives of these bases.

			\item \textbf{Clarity}

			We do not speak of these objects as existing outside of a scalar field.  Further, we can see by these definitions that ``covectors" (how the object interacts with the environment/scalar field) operate on ``vectors" (the concrete measures of the object).  Yet, they do not exist independently.  One describes an element in a vector space and the other describes the vector space using those elements.  Together they form the local geometry (sense of shape, change in direction and curvature).  Indeed, we may even refer to the scope of covectors as the ``covector space" or ``dual space".

		\item \textbf{Changing vectors}

		There are many aspects of this object that can change.  We can translate it, stretch/contract it, and rotating.  Translating it also indicates that if the ambient space that the object exists in also changes shape, for example, moves on the surface of a ball, then that alone would change the values of magnitude and direction without stretching/contracting or rotating.  Thus, taking the derivative (analyzing how it changes locally) will depend on both its ``movement" and the subtle change of the surface.
		\end{itemize}

\end{description}
\end{description}

\newpage
\section*{AI Questions}

\begin{description}

\subsection*{Constrasting vector and covector}

\item \textbf{How do covectors considered passing through level surfaces?\footnote{See the YoutTube video ``Tensors for Beginners 4: What are covectors?" which really helps to clarify things.}}

Good instinct—but let me sharpen it, because there’s a subtle mix-up in that phrasing.

\DEFINE{It’s actually covectors (not contravariant vectors) that define level surfaces.}

\textbf{Contravariant vectors are what *pass through* those surfaces}.

Let’s make that precise.

---

\begin{description}
\item \textbf{1. Level surfaces come from covectors}

Take a scalar function $u(x)$. Its level surfaces are:
\begin{align*}
u(x) = \text{constant}
\end{align*}

The differential:
\begin{align*}
du = \partial_i u \, dx^i
\end{align*}
is a \DEFINE{covector}.

\DEFINE{Geometric meaning:}

\begin{itemize}
\item $du$ is \DEFINE{normal to the level surfaces}
\item it defines a family of parallel “sheets”
\end{itemize}

---

\item \textbf{2. What does a vector do?}

Take a vector:
\begin{align*}
v = v^i \partial_i
\end{align*}

It acts on (u):
\begin{align*}
v(u) = v^i \partial_i u = du(v)
\end{align*}

---

\item \textbf{3. The key interpretation}

\begin{align*}
\boxed{
du(v) = \text{rate at which } v \text{ crosses level surfaces of } u
}
\end{align*}

\begin{itemize}
\item If $du(v) > 0$: moving “upward” through surfaces
\item If $du(v) = 0$: tangent to the surface
\item If $du(v) < 0$: moving “downward”
\end{itemize}

---

\item \textbf{4. So who defines what?}

\begin{itemize}
\item \DEFINE{Covector $du$} → defines the level surfaces
\item \DEFINE{Vector $v$} → passes through them
\end{itemize}

---

\item \textbf{5. Visual picture}

Think of:

\begin{itemize}
\item Level surfaces = parallel sheets of paper
\item Covector = orientation + spacing of the sheets
\item Vector = arrow piercing through them
\end{itemize}

---

\item \textbf{6. Why contravariant vectors are “passing through”}

A contravariant vector has components $v^i$, so:
\begin{align*}
du(v) = \partial_i u \, v^i
\end{align*}

This contraction measures:

\begin{itemize}
\item how aligned the vector is with the normal direction
\item how quickly it crosses the surfaces
\end{itemize}

---

\item \textbf{7. Special cases}

\item \DEFINE{(A) Tangent vector}

\begin{align*}
du(v) = 0
\end{align*}

The vector lies entirely \DEFINE{within} the surface.

---

\item \DEFINE{(B) Normal direction}

If:
\begin{align*}
v^i \propto g^{ij}\partial_j u
\end{align*}

then $v$ points normal to the surface.

---

\item \textbf{8. Why this is inherently covariant vs contravariant}

The pairing:
\begin{align*}
du(v) = \partial_i u \, v^i
\end{align*}

works because:

\begin{itemize}
\item $du$ is covariant (lower index)
\item $v$ is contravariant (upper index)
\end{itemize}

\DEFINE{This ensures the result is coordinate-invariant.}


---

\item \textbf{9. Clean conceptual summary}

\begin{align*}
\boxed{
\text{Covectors define level surfaces; vectors measure how you move through them.}
}
\end{align*}

---

\item \textbf{10. Even sharper intuition}

\begin{itemize}
\item Covector = “stack of surfaces”
\item Vector = “direction of motion”
\item Pairing = “how fast you cross the stack”
\end{itemize}

---

\item \textbf{11. Why this matters}

This viewpoint explains:

\begin{itemize}
\item why gradients are covectors
\item why directional derivatives are $du(v$)
\item why normal vectors come from raising indices
\end{itemize}

---
\end{description}

If you want, I can:

\begin{itemize}
\item show how the **gradient vector** arises from $du$ using the metric, or
\item draw the connection to flux and divergence (this idea shows up there beautifully).
\end{itemize}


\subsection*{Non-affine example}
\item \textbf{How are covariant/contravariant concepts effected by non-affine transormation?}

Perfect—polar coordinates are *the* example where everything (covariant vs contravariant, Christoffel symbols, non-affine effects) becomes concrete.

We’ll go step by step and keep everything explicit.

---

\begin{description}
\item \textbf{1. The coordinate transformation}

From Cartesian $(x,y)$ to polar $(r,\theta)$:
\begin{align*}
x = r\cos\theta,\quad y = r\sin\theta
\end{align*}

This is \DEFINE{non-affine} because second derivatives are nonzero.

---

\item \textbf{2. Basis vectors (very important)}

In polar coordinates, the coordinate basis vectors are:
\begin{align*}
\partial_r, \quad \partial_\theta
\end{align*}

But unlike Cartesian ones, they are \DEFINE{not constant}.

In fact:
\begin{itemize}
\item $\partial_r$ points radially outward (i.e., changes direction with each point)
\item $\partial_\theta$ points tangentially (and its length scales with $r$) (changes size as you go further from the origin.  Also, depends on $r$ so changes direction, too.).
\end{itemize}

---

\item \textbf{3. The metric}

Compute:
\begin{align*}
ds^2 = dx^2 + dy^2 = dr^2 + r^2 d\theta^2
\end{align*} So:
\begin{align*}
g_{ij} =
\begin{pmatrix}
1 & 0 \\
0 & r^2
\end{pmatrix}
\end{align*} and
\begin{align*}
g^{ij} =
\begin{pmatrix}
1 & 0 \\
0 & \frac{1}{r^2}
\end{pmatrix}
\end{align*}

---

\item \textbf{4. Covariant vs contravariant components}

\textbf{Let a vector be:}
\begin{align*}
v = v^r \partial_r + v^\theta \partial_\theta
\end{align*}
\textbf{Lowering indices:}
\begin{align*}
v &= v_r \partial_r + v_\theta \partial_\theta \\
v_r &= v^r, \quad v_\theta = r^2 v^\theta 
\end{align*}
\DEFINE{Already you see:}
\begin{itemize}
\item covariant and contravariant components are \DEFINE{different}
\item this difference depends on position $r$
\end{itemize}

---

\item \textbf{5. Christoffel symbols}

Using the metric, we get the nonzero ones:
\begin{align*}
\Gamma^r_{\theta\theta} &= -r \\
\Gamma^\theta_{r\theta} &= \Gamma^\theta_{\theta r} = \frac{1}{r}
\end{align*}
These arise \DEFINE{purely because the coordinates are non-affine}.

---

\item \textbf{6. Covariant derivative of a vector}

Let’s compute $\nabla_\theta v^r$:
\begin{align*}
\nabla_\theta v^r = \partial_\theta v^r + \Gamma^r_{\theta\theta} v^\theta
\end{align*}
So:
\begin{align*}
\boxed{
\nabla_\theta v^r = \partial_\theta v^r - r v^\theta
}
\end{align*}

---

\textbf{Now compare with ordinary derivative}

\begin{align*}
\partial_\theta v^r
\end{align*}

This misses the extra term $-r v^\theta$.

That missing term comes from:

\DEFINE{the fact that the basis vectors change as you move in $\theta$}

---

\item \textbf{7. Covariant derivative of a covector}

Now take a covector:
\begin{align*}
\omega = \omega_r dr + \omega_\theta d\theta
\end{align*}
Compute:
\begin{align*}
\nabla_\theta \omega_r = \partial_\theta \omega_r - \Gamma^k_{\theta r} \omega_k
\end{align*}
Only one term survives:
\begin{align*}
\Gamma^\theta_{\theta r} = \frac{1}{r}
\end{align*}
So:
\begin{align*}
\boxed{
\nabla_\theta \omega_r = \partial_\theta \omega_r - \frac{1}{r},\omega_\theta
}
\end{align*}

---

\item \textbf{8. Compare the two cases}

\begin{tabular}{|c|c|}
\hline
 Object               & Covariant derivative  \\ 
\hline
 Vector $v^r$         & $+\Gamma v$           \\
 Covector $\omega_r$  & $-\Gamma \omega$      \\
\hline
\end{tabular}

\DEFINE{The \DEFINE{sign difference} is now very concrete.}

---

\item \textbf{9. Geometric meaning (this is the key insight)}

Move around a circle (change $\theta$):

\begin{itemize}
\item The radial direction $\partial_r$ \DEFINE{rotates}
\item The angular direction $\partial_\theta$ \DEFINE{changes length}
\end{itemize}

So even if a vector’s components stay constant, the vector itself changes.

\DEFINE{The Christoffel terms correct for this.}

---

\item \textbf{10. A very concrete example}

Take the vector field:
\begin{align*}
v = \partial_\theta
\quad \Rightarrow \quad v^r = 0,; v^\theta = 1
\end{align*}
Compute:
\begin{align*}
\nabla_\theta v^r = -r
\end{align*}
So:
\begin{align*}
\nabla_\theta v \neq 0
\end{align*}

Even though the components are constant!

---

\textbf{Interpretation}

\begin{itemize}
\item You are moving tangentially around a circle
\item The direction $\partial_\theta$ keeps rotating
\item So the vector is \DEFINE{not constant geometrically}
\end{itemize}

---

\item \textbf{11. Why this only happens in non-affine coordinates}

In Cartesian coordinates:

\begin{itemize}
\item basis vectors don’t change
\item Christoffel symbols = 0
\item ordinary derivatives work
\end{itemize}

In polar coordinates:

\begin{itemize}
\item basis vectors change with position
\item Christoffel symbols $\ne$ 0
\item need covariant derivatives
\end{itemize}

---

\item \textbf{12. Final takeaway}

\begin{align*}
\boxed{
\text{Non-affine transformations make the basis vary with position}
}
\end{align*}

This causes:

\begin{itemize}
\item Christoffel symbols to appear
\item Covariant vs contravariant distinction to matter
\item Ordinary derivatives to fail
\item Covariant derivatives to fix everything
\end{itemize}

---

\item \textbf{13. One-sentence intuition}

\begin{align*}
\boxed{
\text{In polar coordinates, even “constant components” don’t mean a constant vector.}
}
\end{align*}

---
\end{description}

If you want, I can next:

\begin{itemize}
\item compute the \DEFINE{geodesic equations} in polar coordinates (you’ll see circular motion emerge), or
\item show how the \DEFINE{Riemann tensor is still zero} despite nonzero Christoffel symbols (very illuminating).
\end{itemize}


\subsection*{Christoffel Symbols.}
\item \textbf{Help me understand more about Christoffel Symbols.}

The \DEFINE{Christoffel symbols} look mysterious at first, but they’re really just bookkeeping for how your coordinate system bends and twists.

Let’s build intuition step by step.

---

\item \textbf{ 1. What problem are they solving?}

In flat $\mathbb{R}^n$, derivatives are easy:
\begin{align*}
\partial_i V^j
\end{align*}
works fine because the coordinate axes don’t change.

But on a \DEFINE{curved manifold} (or even curvilinear coordinates like polar coordinates), the \DEFINE{basis vectors themselves vary with position}.

So when you differentiate a vector field, two things change:

\begin{enumerate}
\item the components $V^j$
\item the basis vectors $e_j$
\end{enumerate}

Christoffel symbols correct for (2).

---

\item \textbf{ 2. Definition}

The Christoffel symbols $\Gamma^k_{ij}$ are defined by:
\begin{align*}
\nabla_i e_j = \Gamma^k_{ij} e_k.
\end{align*}

Interpretation:

\begin{itemize}
\item  Take the basis vector $e_j$
\item  Differentiate it in direction $x^i$
\item  The result is another vector, expressed as a linear combination of the basis
\end{itemize}

So $\Gamma^k_{ij}$ tells you:

\DEFINE{how the $j\ITH$ basis vector rotates toward the $k\ITH$ direction when moving in the $i\ITH$ direction.}

---

\item \textbf{ 3. Covariant derivative of a vector}

Using this, the correct derivative of a vector field is:
\begin{align*}
\nabla_i V^k = \partial_i V^k + \Gamma^k_{ij} V^j.
\end{align*}

Two parts:

\begin{itemize}
\item $\partial_i V^k$: change in components
\item  $\Gamma^k_{ij} V^j$: correction for moving basis
\end{itemize}

---

\item \textbf{ 4. Why they are NOT tensors}

Christoffel symbols are \DEFINE{not tensors} because:

\begin{itemize}
\item they depend on the coordinate system
\item you can make them vanish at a point by choosing normal coordinates
\end{itemize}

Intuition:
They describe the coordinate system as much as the geometry.

---

\item \textbf{ 5. Formula in terms of the metric}

If you have a Riemannian metric $g_{ij}$, then:
\begin{align*}
\Gamma^k_{ij} = \frac{1}{2} g^{k\ell} \left( \partial_i g_{j\ell} + \partial_j g_{i\ell} - \partial_\ell g_{ij} \right).
\end{align*}

 This shows:

\begin{itemize}
\item Christoffel symbols come entirely from how the \DEFINE{metric changes}
\item No curvature yet — just first derivatives of the metric
\end{itemize}

---

\item \textbf{ 6. Geometric intuition}

Think of walking on Earth:

\begin{itemize}
\item Moving north changes what you mean by “east”
\item So even if a vector “looks constant,” it actually rotates
\end{itemize}

Christoffel symbols measure that rotation.

---

\item \textbf{ 7. Simple example: polar coordinates}

In $(r,\theta)$, the metric is:
\begin{align*}
ds^2 = dr^2 + r^2 d\theta^2.
\end{align*}

Some Christoffel symbols are:
\begin{align*}
\Gamma^r_{\theta\theta} = -r, \quad \Gamma^\theta_{r\theta} = \frac{1}{r}.
\end{align*}

Meaning:

\begin{itemize}
\item moving in $\theta$ changes the radial direction
\item the angular basis vector shrinks/expands with $r$
\end{itemize}

---

\item \textbf{ 8. Connection to curvature}

Christoffel symbols themselves are not curvature, but:

\DEFINE{curvature = derivatives of Christoffel symbols + quadratic terms}

So they are like the “first layer” of geometry.

---

\item \textbf{ 9. Connection to your earlier question (normal derivatives)}

When you computed:
\begin{align*}
\nabla_i \nabla_j u
\end{align*}

you were using Christoffel symbols implicitly:
\begin{align*}
\nabla_i \nabla_j u = \partial_i \partial_j u - \Gamma^k_{ij} \partial_k u.
\end{align*}

 That extra term is exactly the same phenomenon:

\DEFINE{derivatives must account for changing geometry.}

---

\item \textbf{ 10. Mental model (most important)}

A very useful way to think of Christoffel symbols:

\DEFINE{They measure how much your coordinate grid is ``curving” or ``tilting” as you move.}

\begin{itemize}
\item If coordinates are straight (Cartesian): $\Gamma = 0$
\item If coordinates bend (polar, sphere, manifold): $\Gamma \neq 0$

---
\end{itemize}

\HLINE
\newpage
\subsection*{Reimann-Christoffel tensor.}
\item \textbf{Help me understand the Reimann-Christoffel tensor.}

The \DEFINE{Riemann–Christoffel tensor} (usually just called the \DEFINE{Riemann curvature tensor}) is what finally measures \textbf{true curvature}, not just coordinate effects. Let’s build it from the ground up in the clearest way possible.

---

\begin{description}

\item \textbf{1. What is it trying to measure?}

Christoffel symbols tell you how basis vectors change:
\begin{align*}
\nabla_i e_j = \Gamma^k_{ij} e_k.
\end{align*}

But Christoffel symbols can vanish at a point just by changing coordinates — so they \DEFINE{don’t measure real curvature}.

The Riemann tensor answers a deeper question:

\textbf{ If you move a vector around a tiny loop, does it come back unchanged?}

\begin{itemize}
\item Flat space → it comes back the same
\item Curved space → it comes back rotated
\end{itemize}

That rotation is exactly what the Riemann tensor measures.

---

\item \textbf{2. Definition using covariant derivatives}

The Riemann tensor is defined by
\begin{align*}
(\nabla_i \nabla_j - \nabla_j \nabla_i)V^k = R^k_{\ \ell ij} V^\ell.
\end{align*}

So it measures the \DEFINE{failure of second covariant derivatives to commute}.

---

\item \textbf{3. Concrete formula}

In coordinates, the Riemann–Christoffel tensor is

\begin{align*}
R^k_{\ \ell ij} = \partial_i \Gamma^k_{j\ell} - \partial_j \Gamma^k_{i\ell} + \Gamma^k_{im}\Gamma^m_{j\ell} - \Gamma^k_{jm}\Gamma^m_{i\ell}.
\end{align*}

---

\item \textbf{4. What each term means}

Think of the formula in two parts:

\begin{itemize}
\item (A) Derivatives of Christoffel symbols

\begin{align*}
\partial_i \Gamma^k_{j\ell} - \partial_j \Gamma^k_{i\ell}
\end{align*}

This measures whether the “rotation of the basis” changes from point to point.

\item (B) Quadratic correction terms

\begin{align*}
\Gamma^k_{im}\Gamma^m_{j\ell} - \Gamma^k_{jm}\Gamma^m_{i\ell}
\end{align*}

These appear because the basis is already rotating while you move.
\end{itemize}

---

\item \textbf{5. Why this measures curvature}

Imagine moving a vector in two steps:

\begin{enumerate}
\item Move in direction $x^i$, then in direction $x^j$
\item Move in direction $x^j$, then in direction $x^i$
\end{enumerate}

In flat space → same result
In curved space → not the same

The difference equals
\begin{align*}
R^k_{\ \ell ij}V^\ell.
\end{align*}

So the Riemann tensor measures how much the vector \DEFINE{twists due to curvature}.

---

\item \textbf{6. Geometric intuition}

A very concrete picture:

\textbf{Walk on a sphere}

Start at the equator:

\begin{itemize}
\item Walk north
\item Walk east
\item Walk south
\end{itemize}

Your direction changes even though you never turned.

That change in direction = curvature = Riemann tensor.

---

\item \textbf{7. What the indices mean $R^k_{\ \ell ij}$}

Each index has a meaning:

\begin{itemize}
\item $i, j$: directions you move in
\item $\ell$: direction of the vector you are transporting
\item $k$: direction it ends up pointing
\end{itemize}

So the tensor answers:

\DEFINE{ If I transport a vector in direction $\ell$ around the $i$-$j$ plane, how much does it rotate toward direction $k$?}

---

\item \textbf{8. Why it’s a true tensor}

Unlike Christoffel symbols:

\begin{itemize}
\item You cannot make the Riemann tensor vanish by changing coordinates (unless the space is actually flat).
\item If $R^k_{\ \ell ij} = 0$ everywhere → the space is flat.
\end{itemize}

---

\item \textbf{9. Important special cases}

In 2D surfaces

The Riemann tensor contains only \DEFINE{one independent number}:
\begin{align*}
K = \text{Gaussian curvature}.
\end{align*}

So for surfaces:

\begin{itemize}
\item sphere → $K>0$
\item plane → $K=0$
\item saddle → $K<0$
\end{itemize}

---

\item \textbf{10. Most useful mental picture}

Think of the hierarchy:

\begin{tabular}{|l|l|}
\hline
 Object                              & What it measures     \\
\hline
 Metric $g_{ij}$                     & distances            \\
 Christoffel symbols $\Gamma^k_{ij}$ & how coordinates bend \\
 Riemann tensor $R^k_{\ \ell ij}$    & real curvature       \\
\hline
\end{tabular}
---

\newpage
\subsection*{Curvature Explained}
\item \textbf{Evaluate the curvature tensor $B^{\alpha\beta}$ from its definiion.}

You’re now right at the interface between ambient geometry and surface geometry—this is exactly where the \DEFINE{second fundamental form} (curvature tensor of the surface) appears.

Let’s go carefully and clean up the notation as we go.

---

\begin{description}
\item \textbf{1. Setup and notation}

\begin{itemize}
\item $\mathbf{Z}^i$: ambient basis (in $\mathbb{R}^3$ or a manifold)
\item $\mathbf{S}_\alpha = \partial_\alpha \mathbf{x}$: tangent basis to the surface
\item $N_i$: components of the unit normal
\end{itemize}

The \DEFINE{shift tensor} is:
\begin{align*}
Z^i_{\ \alpha} = \mathbf{Z}^i \cdot \mathbf{S}_\alpha
\end{align*}

$\implies$ This converts \DEFINE{surface indices} $\iff$ \DEFINE{ambient indices}.

---

\item \textbf{2. What is $B_{\alpha\beta}$?}

You defined:
\begin{align*}
B_{\alpha\beta} = N_i \nabla_\alpha Z^i_{ \beta}
\end{align*}

This is exactly the \DEFINE{second fundamental form}.

---

\item \textbf{3. Expand the covariant derivative}

\begin{align*}
\nabla_\alpha Z^i_{\ \beta} = \partial_\alpha Z^i_{\ \beta} + \Gamma^i_{\ jk} Z^j_{\ \alpha} Z^k_{\ \beta} - \Gamma^\gamma_{\alpha\beta} Z^i_{\ \gamma}
\end{align*}

So:
\begin{align*}
B_{\alpha\beta} = N_i\left( \partial_\alpha Z^i_{\ \beta} + \Gamma^i_{\ jk} Z^j_{\ \alpha} Z^k_{\ \beta} - \Gamma^\gamma_{\alpha\beta} Z^i_{\ \gamma} \right)
\end{align*}

---

\item \textbf{4. Key simplification}

Since $N_i Z^i_{\ \gamma} = 0$ (normal is orthogonal to tangent vectors), the last term vanishes:

\begin{align*}
N_i \Gamma^\gamma_{\alpha\beta} Z^i_{\ \gamma} = 0
\end{align*}

So:
\begin{align*}
B_{\alpha\beta} = N_i \left( \partial_\alpha Z^i_{\ \beta} + \Gamma^i_{\ jk} Z^j_{\ \alpha} Z^k_{\ \beta} \right)
\end{align*}

---

\item \textbf{5. Recognize ambient covariant derivative}

The expression in parentheses is just:
\begin{align*}
\nabla_\alpha Z^i_{\ \beta} = \nabla_\alpha \mathbf{S}_\beta
\end{align*}

So:
\begin{align*}
\boxed{
B_{\alpha\beta} = N \cdot \nabla_\alpha \mathbf{S}_\beta }
\end{align*}

---

\item \textbf{6. Final geometric form}

\begin{align*}
\boxed{
B_{\alpha\beta} = \mathbf{N} \cdot \frac{\partial^2 \mathbf{x}}{\partial u^\alpha \partial u^\beta} }
\end{align*}

(in Euclidean space, where ambient Christoffel symbols vanish)

---

\item \textbf{7. Interpretation}

\begin{align*}
\boxed{
B_{\alpha\beta}
\text{ measures how the surface bends in the normal direction}
}
\end{align*}

\begin{itemize}
\item It takes two tangent directions $\alpha,\beta$
\item Returns the \DEFINE{normal component of the second derivative}
\end{itemize}

---

\item \textbf{8. Equivalent forms}

You’ll often see:

\textbf{(A) Shape operator form}

\begin{align*}
B_{\alpha\beta} = - \mathbf{S}_\beta \cdot \nabla_\alpha \mathbf{N}
\end{align*}
  

\textbf{ (B) Weingarten map}

\begin{align*}
\nabla_\alpha \mathbf{N} = -B_\alpha^{\beta} \mathbf{S}_\beta
\end{align*}

---

\item \textbf{9. Key properties}

\begin{align*}
\boxed{
B_{\alpha\beta} = B_{\beta\alpha}
}
\end{align*}

\textbf{symetric}.

---

\item \textbf{10. Big picture}

\begin{align*}
\boxed{
\text{First fundamental form } g_{\alpha\beta} \rightarrow \text{intrinsic geometry}
}
\end{align*}
\begin{align*}
\boxed{
\text{Second fundamental form } B_{\alpha\beta} \rightarrow \text{extrinsic curvature}
}
\end{align*}

---

\item \textbf{11. One-sentence intuition}

\begin{align*}
\boxed{
B_{\alpha\beta} = \text{“how the tangent plane rotates as you move along the surface”}
}
\end{align*}

---
\end{description}

If you want, I can:

\begin{itemize}
\item compute $B_{\alpha\beta}$ explicitly for a sphere (you’ll see it equals the metric!), or
\item derive Gaussian curvature $K = \det(B)/\det(g)$, which ties everything together beautifully.
\end{itemize}


\end{description}

\subsection*{Relationship between $B_{\alpha\beta}$ and $B^\alpha_\beta$}
\item \textbf{$B_{\alpha\beta} \to B^\alpha_\beta$}

This is a really important distinction—you're moving from \DEFINE{geometry} to \DEFINE{linear algebra on the tangent space}.

---

\begin{description}
\item \textbf{1. What you already have}

The object
\begin{align*}
B_{\alpha\beta}
\end{align*}
is the \DEFINE{second fundamental form}:
\begin{align*}
B_{\alpha\beta} = N \cdot \nabla_\alpha \mathbf{S}_\beta
\end{align*}

$\implies$ It takes \DEFINE{two tangent directions} and gives a scalar.

So:
\begin{align*}
\boxed{
B_{\alpha\beta} \text{ is a } (0,2)\text{-tensor (bilinear form)}
}
\end{align*}

---

\item \textbf{2. Raising an index}

Define:
\begin{align*}
B^\alpha_{\ \beta} = g^{\alpha\gamma} B_{\gamma\beta}
\end{align*}

$\implies$ You used the metric to raise one index.

---

\item \textbf{3. What changes conceptually?}

\begin{description}
\item (A) $B_{\alpha\beta}$: bilinear form

\begin{align*}
B_{\alpha\beta} X^\alpha Y^\beta
\end{align*}

\begin{itemize}
\item \textbf{Input: two vectors}
\item \textbf{Output: a number}
\item \textbf{Measures \DEFINE{normal curvature}}
\end{itemize}

---

\item (B) $B^\alpha_{\ \beta}$: linear operator

\begin{align*}
(B^\alpha_{\ \beta} V^\beta)
\end{align*}

\begin{itemize}
\item \textbf{Input: one vector}
\item \textbf{Output: another vector}
\end{itemize}
\end{description}

So:
\begin{align*}
\boxed{
B^\alpha_{\ \beta} : T_pS \to T_pS
}
\end{align*}

This is the \DEFINE{shape operator (Weingarten map)}.

---

\item \textbf{4. Geometric meaning of the difference}

\item \textbf{$B_{\alpha\beta}$}

\begin{itemize}
\item \textbf{Measures \DEFINE{how much the surface bends}}
\item \textbf{Think: “curvature as a quadratic form”}
\end{itemize}

---

\item \textbf{$B^\alpha_{\ \beta}$}

\begin{itemize}
\item \textbf{Tells you \DEFINE{how tangent vectors are rotated}}
\item \textbf{Think: “curvature as a transformation”}
\end{itemize}

---

\item \textbf{5. The key identity connecting them}

\begin{align*}
\boxed{
\nabla_\alpha \mathbf{N} = -\textbf{B}^\beta_{\alpha} \mathbf{S}_\beta
  }
\end{align*}

$\implies$ This shows:

\item \textbf{$B^\alpha_{\ \beta}$ governs how the \DEFINE{normal changes}}
\item \textbf{That’s why it’s an \DEFINE{operator}}

---

\item \textbf{6. Eigenvalues = principal curvatures}

This is the big payoff:

\begin{align*}
B^\alpha_{\ \beta} v^\beta = \kappa v^\alpha
\end{align*}

\begin{itemize}
\item \textbf{Eigenvalues $\kappa_1, \kappa_2$ = \DEFINE{principal curvatures}}
\item \textbf{Eigenvectors = \DEFINE{principal directions}}
\end{itemize}

---

\item \textbf{7. Why you can’t do this with $B_{\alpha\beta}$}

You \textit{cannot} take eigenvalues of $B_{\alpha\beta}$ directly because:

\begin{itemize}
\item \textbf{It maps $ (V,W) \to \mathbb{R} $}
\item \textbf{Not $ V \to V $}
\end{itemize}

Raising an index converts it into something you can diagonalize.

---

\item \textbf{8. Matrix analogy}

Think of:

\item \textbf{Metric:}

\begin{align*}
g_{\alpha\beta}
\end{align*}

\item \textbf{Second fundamental form:}

\begin{align*}
B_{\alpha\beta}
\end{align*}

Then:
\begin{align*}
B^\alpha_{\ \beta} = g^{-1} B
\end{align*}

$\implies$ Exactly like:
\begin{align*}
A = G^{-1}B
\end{align*}

which turns a quadratic form into a linear map.

---

\item \textbf{9. One clean summary}

\begin{align*}
\boxed{
\begin{aligned}
B_{\alpha\beta} &:\ \text{measures curvature} \\
B^\alpha_{\ \beta} &:\ \text{acts as curvature operator}
\end{aligned}
}
\end{align*}

---

\item \textbf{10. Intuition}

* $B_{\alpha\beta}$: “how curved is the surface in these two directions?” \\
* $B^\alpha_{\ \beta}$: “if I move this vector, how does the surface bend it?”

---
\end{description}

If you want, I can:

* compute both explicitly for a sphere (you’ll see $B^\alpha_{\ \beta} = \delta^\alpha_\beta$), or \\
* show how Gaussian curvature $K = \det(B^\alpha_{\ \beta})$ drops out immediately.

\end{description}

\HLINE
\newpage
\subsection*{Describe the Calculus of Variations}

Calculus of variations is about a different kind of optimization:

\begin{align*}
\boxed{
\text{Instead of minimizing numbers, you minimize functions (or curves, surfaces, fields).}
}
\end{align*}

---

\begin{description}
\item \textbf{1. The basic problem}

In ordinary calculus:
\begin{align*}
\text{minimize } f(x)
\end{align*}

In calculus of variations:
\begin{align*}
\boxed{
\text{minimize } J[u] = \int_a^b L(x, u(x), u'(x))\,dx
}
\end{align*}

\begin{itemize}
\item $u(x)$: unknown \DEFINE{function}
\item $J[u]$: a \DEFINE{functional} (a function of a function)
\item $L$: called the \DEFINE{Lagrangian}
\end{itemize}

---

\item \textbf{2. What does ``varying a function" mean?}

We don’t change $x$—we change the whole function:

\begin{align*}
u(x) \to u(x) + \varepsilon \eta(x)
\end{align*}

\begin{itemize}
\item $\eta(x$): arbitrary “test function”
\item $\varepsilon$: small parameter
\end{itemize}

---

\item \textbf{3. The key idea (like derivatives)}

We require:
\begin{align*}
\boxed{
\frac{d}{d\varepsilon} J[u + \varepsilon \eta]\Big|_{\varepsilon=0} = 0
}
\end{align*}

This is the analogue of:
\begin{align*}
f'(x)=0
\end{align*}

---

\item \textbf{4. Compute the variation}

\begin{align*}
J[u+\varepsilon \eta] = \int_a^b L(x, u+\varepsilon\eta, u'+\varepsilon\eta')\,dx
\end{align*}

Differentiate:
\begin{align*}
\delta J = \int_a^b \left( \frac{\partial L}{\partial u}\eta + \frac{\partial L}{\partial u'}\eta' \right) dx
\end{align*}

---

\item \textbf{5. Integration by parts}

Move derivative off $\eta'$:
\begin{align*}
\int_a^b \frac{\partial L}{\partial u'} \eta'\,dx = \left[ \frac{\partial L}{\partial u'} \eta \right]_a^b - \int_a^b \frac{d}{dx}\left(\frac{\partial L}{\partial u'}\right)\eta\,dx
\end{align*}

Assume $\eta(a)=\eta(b)=0$, so boundary term vanishes.

---

\item \textbf{6. Euler–Lagrange equation}

We get:
\begin{align*}
\delta J = \int_a^b \left( \frac{\partial L}{\partial u} - \frac{d}{dx}\frac{\partial L}{\partial u'} \right)\eta\,dx
\end{align*}

Since $\eta$ is arbitrary:

\begin{align*}
\boxed{
\frac{\partial L}{\partial u} - \frac{d}{dx}\left(\frac{\partial L}{\partial u'}\right) = 0
}
\end{align*}

---

\item \textbf{7. This is the main result}
\begin{align*}
\boxed{
\text{Extrema of functionals satisfy the Euler–Lagrange equation}
}
\end{align*}

---

\item \textbf{8. Example: shortest path (geodesics)}

Minimize length:
\begin{align*}
J[y] = \int_a^b \sqrt{1 + (y')^2}\,dx
\end{align*}

Then:
\begin{align*}
L = \sqrt{1+(y')^2}
\end{align*}

Euler–Lagrange gives:
\begin{align*}
\frac{d}{dx}\left(\frac{y'}{\sqrt{1+(y')^2}}\right)=0
\end{align*}

Solution:
\begin{align*}
y' = \text{constant} \Rightarrow y = ax + b
\end{align*}

\textbf{Straight lines are shortest paths.}

---

\item \textbf{9. Physical meaning}

Calculus of variations underlies:

\begin{description}
\item  \DEFINE{Mechanics}

\begin{align*}
\text{Action } S = \int L\,dt
\end{align*}

\begin{align*}
\boxed{
\delta S = 0 \Rightarrow \text{Newton’s laws}
}
\end{align*}

---

\item  \DEFINE{General relativity}

\begin{itemize}
\item Geodesics = extremize spacetime interval
\end{itemize}

---

\item  \DEFINE{Field theory}

\begin{itemize}
\item Maxwell’s equations
\item Einstein’s equations
\end{itemize}

All come from variational principles.

\end{description}
---

\item \textbf{10. Higher-dimensional version}

For fields $u(x^\mu$):

\begin{align*}
J[u] = \int L(x^\mu, u, \partial_\mu u)\, d^n x
\end{align*}

Euler–Lagrange becomes:
\begin{align*}
\boxed{
\frac{\partial L}{\partial u} - \partial_\mu \left(\frac{\partial L}{\partial (\partial_\mu u)}\right) = 0 }
\end{align*}

---

\item \textbf{11. Big picture}

\begin{align*}
\boxed{
\text{Calculus of variations turns optimization problems into differential equations}
}
\end{align*}

---

\item \textbf{12. Intuition}

\begin{itemize}
\item You “wiggle” a function slightly
\item If it’s optimal, no wiggle lowers the value
\item That condition forces a differential equation
\end{itemize}

---

\item \textbf{13. One-sentence summary}

\begin{align*}
\boxed{
\text{Find the function whose small variations don’t change the integral.}
}
\end{align*}
\end{description}

---

If you want, I can:

\begin{itemize}
\item derive geodesics on a sphere using this, or
\item connect this directly to your earlier curvature tensors (they appear naturally in variational problems).
\end{itemize}


\end{document}
